\begin{center}\small\begin{tabular}{ll}\hline Illustrative input & Value\\\hline
Mass $m$; radius $R$ & $1\,{\rm kg}$; $0.05\,{\rm m}$\\
Inertia $I$; lengths $\ell,\ell_c$ & $0.001\,{\rm kg\,m^2}$; $0.05,0.05\,{\rm m}$\\
Incoming normal velocity; axial spin & $-1\,{\rm m/s}$; $20\,{\rm rad/s}$\\
Prescribed $e_n$; free axial impulse $k_n$ & $0.6$; $-0.008\,{\rm kg\,m^2/s}$\\\hline\end{tabular}\end{center}
The torque impulse is prescribed solely to demonstrate the missing degree of freedom; it is not inferred from a measured outcome or claimed as a material prediction.
\begin{center}\small\begin{tabular}{lrr}\hline Computed quantity & Force impulse only & Force plus torque impulse\\\hline
Outgoing axial spin, rad/s & 20.000 & 12.000\\
Outgoing axial angular momentum, kg m$^2$/s & 0.020 & 0.012\\
Outgoing normal velocity, m/s & 0.600 & 0.600\\
\hline\end{tabular}\end{center}
The full-wrench energy changes from $0.700$ to $0.252\,{\rm J}$. The axial angular-momentum change is exactly the prescribed torque impulse, while the central normal force impulse contributes no angular impulse.
